Zusammenfassend lässt sich festhalten, dass im Rahmen dieser Arbeit erfolgreich eine immersive Virtual-Reality-Anwendung zur Simulation der Fortbewegung im Rollstuhl entwickelt wurde. Das Ziel bestand darin, Barrieren im urbanen und baulichen Umfeld erfahrbar zu machen und das Bewusstsein für die Herausforderungen von Rollstuhlfahrenden zu stärken.

Die Ergebnisse zeigen, dass insbesondere die Kombination aus virtueller Umgebung und physischer Simulation durch den entwickelten Hardware-Simulator ein intensives und realistisches Nutzererlebnis ermöglicht. Die Integration von haptischem Feedback, wie Neigung und Vibrationen, trägt wesentlich zur Immersion bei und hebt die Anwendung von herkömmlichen VR-Systemen ab.

Zukünftige Arbeiten könnten sich auf die Weiterentwicklung der Hardware, die Erweiterung der simulierten Szenarien sowie die Durchführung umfangreicher Nutzerstudien zur Evaluation der Wirkung konzentrieren.
Darüber hinaus empfiehlt es sich, die virtuelle Realität um eine größere Vielfalt an Hindernissen und Interaktionselementen zu erweitern, um eine noch realitätsnähere und anspruchsvollere Umgebung zu schaffen. In diesem Zusammenhang sollten diese zusätzlichen Hindernisse auch in die Hardware-Simulation integriert werden, sodass sie nicht nur visuell, sondern ebenfalls physisch korrekt erfahrbar sind und demnach eine konsistente Abbildung der virtuellen Umgebung gewährleistet wird.